\chapter{Dove guardare}
\label{cap:dove_guardare}

La domanda da cui è partito questo libro si è progressivamente trasformata. Non è più solo "perché il Veneto sembra privo di arte rupestre?" ma qualcosa di più preciso: quali tipi di luoghi, in questo territorio, non sono mai stati cercati con i metodi giusti? La risposta non è una lista di coordinate, ma una logica: capire come si muovevano gli animali cacciati, come si muovevano i cacciatori, come funzionano i metodi sistematici sviluppati altrove, e come cambia tutto questo al variare del periodo. Ogni epoca ha una firma diversa nei depositi rupestri, e ogni firma richiede strumenti di ricerca adattati.

% ============================================================
\section{La logica del movimento: animali e uomini}
% ============================================================

Il primo errore da evitare è ragionare per corridoi. L'immagine del cacciatore preistorico che segue la selvaggina lungo vallate continue, risalendo dalle quote invernali alle quote estive secondo percorsi prevedibili, è intuitiva ma empiricamente fragile, almeno per le specie che hanno dominato la fauna alpina veneta dalla fine del Pleistocene in poi.

Il camoscio (\emph{Rupicapra rupicapra}) e lo stambecco alpino (\emph{Capra ibex}) non sono animali di lunga migrazione. Studi radiotelemetrici su popolazioni alpine contemporanee indicano home range che raramente superano i venti chilometri quadrati per il camoscio e i trenta per lo stambecco; in condizioni di densità normale, questi animali non abbandonano la fascia altitudinale che conoscono, e i movimenti stagionali avvengono entro pochi chilometri verticali, non lungo centinaia di chilometri di versante. Questo non significa che le valli non venissero attraversate, ma che l'attraversamento non era il pattern dominante: l'arte rupestre legata alla caccia segue la distribuzione spaziale degli animali, non i desideri del cacciatore di muoversi velocemente.

Il secondo errore simmetrico è immaginare che il cacciatore preistorico si muovesse senza vincoli energetici. Il modello ARC (Accounting for Roughness and Cost) sviluppato da Herman Pontzer e collaboratori \citep{pontzer2016} quantifica il costo metabolico della locomozione su terreno irregolare per mammiferi di taglia molto diversa. Per un essere umano adulto, il costo energetico di un dislivello positivo è circa quattro volte quello dello spostamento in piano a parità di distanza orizzontale; il costo del dislivello negativo è circa il doppio. Su un percorso alpino tipico con cento metri di salita e ottanta di discesa per chilometro orizzontale, questo si traduce in un aumento del costo totale dell'ordine del quaranta per cento rispetto al piano. Le scelte di localizzazione dei siti non seguono la distanza lineare ma il costo del tragitto.

Uniti, questi due vincoli producono una predizione verificabile: i siti si addenseranno in luoghi che soddisfano contemporaneamente la distanza di costo accettabile dalla base (Kompatscher e Hrozny Kompatscher parlano di due ore di marcia, che su terreno alpino corrispondono a distanze assai variabili a seconda del dislivello), la visibilità sulle zone di pascolo degli animali cacciati, e la disponibilità d'acqua a meno di ottanta metri. La combinazione dei tre fattori restringe enormemente il numero di candidati su qualsiasi cartografia topografica.

% ============================================================
\section{Il metodo Kompatscher: nato dal camminare, non dal calcolare}
% ============================================================

Il metodo messo a punto da Karl Kompatscher e Margit Hrozny Kompatscher nel corso di decenni di lavoro sulle Alpi trentine e alto-atesine non nasce da un modello teorico, ma dalla sistematizzazione di un'esperienza accumulata sul campo \citep{kompatscher2007}. I due ricercatori hanno censito oltre trecento siti in Trentino-Alto Adige applicando una griglia di criteri ricavati a posteriori dall'analisi dei siti già trovati, e verificati poi su aree nuove.

Il metodo si regge su due parametri fondamentali. Il primo è il bacino di due ore: un sito di abitazione temporanea, per essere funzionale, deve essere raggiungibile in circa due ore di marcia da un punto d'acqua affidabile e deve permettere di raggiungere, nelle stesse due ore, le zone di caccia principali. L'area entro due ore di marcia su terreno alpino ha forma irregolare, dipendente dal dislivello, e in media copre una superficie di cinque-quindici chilometri quadrati attorno al sito. Il secondo parametro è il numero di direttrici: un sito è preferenziale se da esso partono almeno due percorsi distinti verso zone di caccia o verso valichi, perché riduce l'imprevedibilità dell'esito di una singola battuta.

Da questi due parametri derivano tre tipi di sito. Il campo base (campo di sosta prolungata, spesso con acqua a venti-ottanta metri, morfologia riparata, tracce di combustione) è il nodo della rete; il sito di caccia e avvistamento è un punto panoramico o un valico da cui si controllano i movimenti degli animali; il punto di sosta temporanea è un riparo minimo usato durante la battuta. L'arte rupestre, quando è presente, tende a concentrarsi nei campi base e nei siti di avvistamento, raramente nei punti di sosta breve.

La quota dell'acqua è un filtro molto selettivo. Su terreno calcareo carsificato, le sorgenti emergono a quote precise legate alla stratigrafia; su terreno cristallino, i ruscelli di fusione nivale funzionano stagionalmente. La distanza di venti-ottanta metri non è arbitraria: al di sotto di venti metri l'acqua è troppo vicina per garantire visibilità sufficiente verso le zone di caccia (il sito è in fondo valle), al di sopra di ottanta metri il costo di approvvigionamento diventa penalizzante nelle operazioni quotidiane di un campo stagionale.

Ciò che è utile di questo metodo per il Veneto non è la lista dei siti trentini trovati, ma la logica sottostante: cercare su cartografia i punti che soddisfano contemporaneamente più direttrici di movimento, acqua a distanza ragionevole, riparo morfologico. Questi punti si possono identificare prima di andare sul campo, riducendo il territorio da esplorare da migliaia a decine di ettari.

% ============================================================
\section{Le prospezioni sistematiche nelle Alpi: tre casi di riferimento}
% ============================================================

Il metodo Kompatscher non è un caso isolato. Tre programmi di ricerca condotti in contesti alpini diversi mostrano convergenze significative nella logica e divergenze istruttive nei risultati, e costituiscono insieme un quadro di riferimento per capire cosa può aspettarsi una prospezione sistematica in Veneto.

\subsection*{Silvretta Historica: 230 siti in 550 chilometri quadrati}

Il programma Silvretta Historica, coordinato da Martin Cornelissen e Tobias Reitmaier tra il 2009 e il 2016, ha coperto sistematicamente circa 550 chilometri quadrati di territorio alpino al confine tra Austria e Svizzera, a quote prevalentemente superiori ai duemila metri \citep{cornelissen2016}. Il risultato è stato il censimento di oltre 230 siti, il novanta per cento dei quali sconosciuto prima dell'avvio del programma. Nessuno di essi era stato trovato nelle prospezioni precedenti, nonostante l'area fosse attraversata da itinerari escursionistici e studiata da ricercatori locali.

La ragione non è che i siti fossero nascosti: è che le prospezioni precedenti non erano state progettate per trovarli. Un camminatore che percorre un sentiero vede le rocce a meno di tre metri dalla traccia; un ricercatore che progetta una prospezione sistematica a transetti paralleli di trenta-cinquanta metri copre l'intero territorio. La Silvretta ha prodotto non solo un catalogo, ma la dimostrazione empirica che un'area descritta come "priva di tracce" prima di una prospezione sistematica può rivelarsi densamente frequentata. Per il Veneto, questo è il riferimento metodologico più diretto: l'assenza dalla letteratura non è prova di assenza dal territorio.

\subsection*{Tirolo e Carinzia: le revisioni sistematiche}

In Alto Adige, Tirolo e Carinzia, una serie di ricercatori ha sviluppato programmi di revisione dell'esistente e di prospezione di nuove aree, con approcci parzialmente diversi da Kompatscher ma convergenti nella logica. Franz Mandl in Stiria \citep{mandl2002} e Walter Leitner in Tirolo \citep{leitner2013} hanno documentato centinaia di siti di quota (principalmente caccia, accampamenti stagionali, strutture pastorali di varie epoche) attraverso prospezioni che privilegiano la copertura sistematica rispetto alla verifica di segnalazioni pregresse.

Ciò che emerge da questi lavori è una distribuzione dei siti che conferma la logica di Kompatscher: i siti si addensano intorno a valichi, in spianate protette dal vento nelle vicinanze di sorgenti, e su punti panoramici con visibilità multipla. La variazione rispetto al modello Kompatscher riguarda principalmente l'adattamento a morfologie diverse (rocce cristalline invece di calcareo-dolomitiche, fondovalli più ampi). Per il Veneto, che condivide con queste aree sia morfologie carbonatiche (Prealpi, Dolomiti) sia aree cristalline (Lagorai, Cima d'Asta), entrambi i modelli sono pertinenti.

\subsection*{Scandinavia: RTI e metodi di rilevamento per l'arte rupestre}

La ricerca sull'arte rupestre scandinava, concentrata soprattutto in Norvegia (Alta, sito UNESCO) e Svezia (Nämforsen, Bohuslän), ha sviluppato negli ultimi vent'anni una metodologia di rilevamento che va oltre la semplice prospezione topografica. La tecnica RTI (Reflectance Transformation Imaging) \citep{drap2013}, applicata sistematicamente a pannelli rupestri, permette di rivelare incisioni superficiali invisibili in luce naturale diretta e non rilevabili con la pulizia meccanica. La stessa tecnica è stata applicata con risultati significativi a Creswell Crags (Inghilterra), dove incisioni paleolitiche sono state identificate su pareti che erano state studiate per decenni senza rilevarle.

La fotogrammetria strutturata (Structure from Motion, SfM) permette invece la costruzione di modelli tridimensionali delle superfici rocciose, utile sia per documentare siti già noti sia per identificare morfologie superficiali compatibili con incisioni profonde in condizioni di luce non favorevole. La luce radente artificiale (raking light), nella versione standardizzata sviluppata in ambito scandinavo, permette campagne notturne che rivelano segni non visibili di giorno.

Questi metodi non sostituiscono la prospezione pedestr, ma la integrano: una volta identificata tramite criteri Kompatscher un'area prioritaria, RTI e fotogrammetria permettono di esaminare sistematicamente le superfici rocciose senza dipendere dalle condizioni di illuminazione naturale.

% ============================================================
\section{Il Paleolitico superiore: grotte, rocce e clima}
% ============================================================

Il Paleolitico superiore nelle Alpi venete (approssimativamente da trentacinquemila a undicimilacinquecento anni prima dell'anno zero) è caratterizzato da condizioni ambientali radicalmente diverse da quelle attuali. Le fasi più fredde (Glaciale Massimo, intorno a ventiquattromila anni prima dell'anno zero) vedono il limite delle nevi perenni scendere fino a seicento-ottocento metri sulle Prealpi; la fauna include renna, rinoceronte lanoso, mammut nelle fasi più antiche, e camoscio, stambecco e cervo nelle fasi di interstadio. La tecnologia è epigravettiana (lamelle a dorso, armature microlitiche) nelle fasi tardo-pleistoceniche; la struttura sociale suggerisce gruppi di venti-quaranta individui con ampi range di spostamento nella fase fredda, tendenti a restringersi con il miglioramento climatico.

Per questo periodo, i siti di arte rupestre attesi in Veneto non si trovano in quota: le Alpi erano in gran parte inospitali o inaccessibili. I siti si concentrano nelle valli prealpine inferiori, nelle grotte e ripari tra cento e seicento metri, nella fascia di transizione tra pianura e montagna dove la selvaggina era disponibile e l'accesso relativamente sicuro. Riparo Tagliente, Riparo Mezzena, Grotta di Fumane (Lessinia) sono tutti in questa fascia; Riparo Dalmeri (Valsugana) a 1240 metri è un caso eccezionale legato a una fase di miglioramento climatico tardo-pleistocenico.

Il tipo di arte atteso per questo periodo è quello documentato a Dalmeri: pitture su supporto mobile (ciottoli dipinti), raramente arte parietale, e quando parietale, in cavità con pareti calcaree lavate o in ripari con fondo di accumulo sedimentario. Le grotte della Lessinia non risultano sistematicamente rilevate per arte parietale con metodi RTI; le valli prealpine della Valsugana e dell'Agordino non risultano oggetto di prospezioni dedicate per arte rupestre di questo periodo. I porfidi del Lagorai (quota inadatta per questo periodo) e le dolomie di vetta possono essere esclusi.

La metodologia appropriata per il Paleolitico è diversa da quella per i periodi successivi: non si cercano valichi e selle, ma cavità o ripari con pareti calcaree tra cento e seicento metri, con attenzione particolare alle superfici interne lavate dall'acqua, alle pareti verticali in penombra, e ai depositi sedimentari che possono coprire pannelli dipinti. RTI e fotogrammetria sono qui lo strumento primario di indagine delle superfici, non la prospezione a transetti.

% ============================================================
\section{Il Mesolitico: valichi, passi e logica di Kompatscher}
% ============================================================

Con il riscaldamento postgiaciale (da circa undicimilacinquecento a settemila anni prima dell'anno zero), la fauna cambia e con essa la strategia di caccia. Il cervo, lo stambecco e il camoscio diventano le specie dominanti alle quote montane; la renna scompare progressivamente dal territorio veneto. Le quote accessibili si alzano: intorno a ottomila anni prima dell'anno zero, le Alpi sopra i duemila metri tornano praticabili nella stagione estiva. La tecnologia è geometrica (trapezi, punte a dorso di piccole dimensioni); la struttura sociale è quella del gruppo di cacciatori-raccoglitori con mobilità stagionale, con siti di quota usati nella stagione calda e siti di valle nella stagione fredda.

Questo è il periodo a cui il metodo Kompatscher è più direttamente applicabile. I criteri di ricerca sono ben definiti: valichi e selle con almeno due direttrici di movimento verso zone di caccia, a quote tra millequattrocento e duemilatrecento metri (con il Castelnoviano che tende a preferire versanti esposti a sud per le temperature mattutine), con sorgente o torrentello a venti-ottanta metri, e morfologia riparata almeno su un lato. L'arte rupestre del Mesolitico, quando presente, si trova spesso in corrispondenza di queste strutture: non pitture elaborate, ma coppelle (bicchierini scolpiti), tratti incisi, impronte di mano, segni geomorfici su superfici calcaree subpianeggianti.

Nel Veneto, i candidati applicando questi criteri si distribuiscono lungo l'arco alpino a quote appropriate. I passi del gruppo dolomitico cadorino (Passo Duran, Passo Cereda, Passo di Giau, col dei passi sopra Alleghe) presentano morfologie adeguate e non risultano sistematicamente cercati per arte rupestre. Il Cansiglio, con le sue doline e i pianori carsici, rappresenta una variante morfologica del modello: non un valico, ma un bacino di confluenza di più percorsi con acqua abbondante (Pian Cansiglio ha sorgenti multiple). La ricerca di Ferrara ha documentato la frequentazione mesolitica in quest'area; l'arte rupestre specifica non è stata ancora il target principale.

% ============================================================
\section{Il Neolitico: coppelle, stele e pianori}
% ============================================================

Con il Neolitico (da circa settemila a cinquemila anni prima dell'anno zero) cambia non solo la tecnologia ma l'assetto insediativo. Le comunità agricole occupano la pianura e i fondovalle in modo più stabile; la montagna è frequentata stagionalmente per la pastorizia e la caccia, ma non con la stessa logica di mobilità mesolitica. L'arte rupestre neolitica ha caratteristiche proprie: le coppelle (cavità circolari o ovali incise per percussione sulla superficie rocciosa) diventano ubiquitarie e spesso si trovano su rocce subpianeggianti o debolmente inclinate, non sulle verticali preferite in altri periodi. Le stele menhir compaiono in questo periodo in vari contesti alpini (Trentino, Lombardia, Valle d'Aosta, Alto Adige) e vengono spesso riutilizzate in strutture medievali, dove possono essere recuperate come materiale di spoglio.

La metodologia per il Neolitico si divide quindi tra ricerca di coppelle e ricerca di stele. Le coppelle si cercano su pianori rocciosi calcarei a quote variabili (dalla pianura alle quote di mezza montagna), con attenzione ai punti panoramici e ai crocevia di sentieri storici; la densità attesa in contesti calcarei europei simili è di uno-cinque siti per dieci chilometri quadrati, ma questa è una stima molto rozza, dipendente dalla morfologia locale. Le stele si cercano in murature di edifici storici (chiese rurali, masi, fontane), nelle collezioni locali di materiale di reimpiego, e in archivi fotografici di lavori edilizi del Novecento. Il recupero di stele dalle murature è metodologicamente diverso dalla prospezione rupestre, ma può produrre risultati significativi: in Trentino, decine di stele sono state identificate proprio in questo modo.

In Veneto, le Prealpi calcaree (Lessinia, Piccole Dolomiti, Alpago) presentano morfologie idonee per le coppelle; i fondovalle storicamente abitati (Valpolicella, Vallagarina, Agordino) potrebbero nascondere stele in reimpiego. La ricerca specifica per questi target non risulta sistematicamente condotta.

% ============================================================
\section{L'Eneolitico: valichi come nodi di scambio}
% ============================================================

L'Eneolitico (da circa cinquemila a tremilacinquecento anni prima dell'anno zero) è il periodo delle prime metallurgie e degli scambi a lunga distanza. Ötzi, l'Uomo del Similaun, appartiene a questo periodo (3300-3100 anni prima dell'anno zero); la sua presenza su un passo alpino d'alta quota nel momento della morte illustra il ruolo che i valichi avevano nei movimenti di persone e merci. In questo periodo, le selle e i valichi diventano nodi di un sistema di scambio di rame, ossidiana, ambra e prodotti agricoli; l'arte rupestre eneolitica tende a concentrarsi in corrispondenza di questi nodi, come documentato in Val Camonica (incisioni sui bordi dei valichi laterali) e nel Tirolo meridionale.

Per il Veneto, i valichi delle Piccole Dolomiti (Passo Borcola, Passo Campogrosso), del Pasubio (Bocca di Campiglia, passo Xomo), del Brenta-Adamello (Passo Tremalzo) e delle Dolomiti bellunesi e cadorine (Passo Duran, Passo Cereda, Forcella Valbona) sono candidati naturali per arte rupestre eneolitica. La morfologia eneolitica preferita è diversa da quella mesolitica: non il campo base riparato di un cacciatore, ma la parete verticale o inclinata in corrispondenza di un passaggio, dove si incidono segni visibili da chi transita. Le superfici calcaree lisce, esposte in verticale o con inclinazione di quaranta-settanta gradi, nel raggio di cinquanta metri da un punto di passaggio obbligato, sono il target specifico per questo periodo.

% ============================================================
\section{L'età del bronzo: pareti verticali e bacini visivi}
% ============================================================

L'età del bronzo (da circa tremilacinquecento a duemilanovecento anni prima dell'anno zero) è il periodo di maggiore intensità documentata per l'arte rupestre in area alpina. Valcamonica e Valtellina hanno migliaia di figure di questo periodo; la Val d'Assa veneta è in larga parte bronzina; Monte Luppia (Garda) ha pannelli riferibili a questo periodo. La firma stilistica è riconoscibile: figure geometriche complesse (retticolati, spirali, cerchi concentrici), figure umane stilizzate (oranti con braccia alzate), armi (asce, pugnali, lance), animali (bovini, cervi) su superfici verticali o leggermente inclinate con bacino visivo ampio.

Il target morfologico per l'età del bronzo è specifico: superficie calcarea verticale o inclinata tra i trenta e i settanta gradi, con overhang (sporgenza soprastante) che protegga dall'erosione pluviale, a quote tra quattrocento e milleduecento metri, con ampia visibilità su almeno una vallata principale. La combinazione di inclinazione, protezione e quota restringe i candidati su qualsiasi cartografia topografica. La metodologia di rilevamento appropriata è la combinazione di fotogrammetria SfM (per il modello 3D della superficie), luce radente artificiale nelle ore notturne (per rivelare le incisioni più superficiali), e RTI su campioni selezionati.

In Veneto, le aree candidate per arte rupestre dell'età del bronzo non ancora sistematicamente cercate includono la Lessinia settentrionale (zona Sant'Anna d'Alfaedo, Vajo dell'Anguilla, pareti calcaree al limite del bosco tra seicento e novecento metri), l'Altopiano di Asiago nelle zone non interessate dai combattimenti della Prima guerra mondiale (parte meridionale, verso Roana e Foza), e il versante meridionale del Grappa verso Romano d'Ezzelino e Pove del Grappa. Tutte e tre le aree hanno morfologie calcaree con pareti verticali nella fascia di quota appropriata; nessuna delle tre risulta oggetto di prospezioni sistematiche dedicate all'arte rupestre.

% ============================================================
\section{L'età del ferro e il periodo venetico: continuità e aggiunte}
% ============================================================

Nell'età del ferro (da circa duemilanovecento a duemila anni prima dell'anno zero) e nel periodo venetico (fino alla romanizzazione), l'arte rupestre non cessa ma spesso si sovrappone ai pannelli bronzini già esistenti. Le aggiunte sono riconoscibili per stile (figure più schematiche, iscrizioni, simboli geometrici elementari) e tecnica (incisione più fine rispetto alla martellina del Bronzo). In Val Camonica, le aggiunte di età del ferro e periodo celtico o romano sono documentate su decine di pannelli già utilizzati in precedenza: il luogo continua ad essere frequentato e segnato, ma con un'intensità e un vocabolario diversi.

La metodologia per identificare arte rupestre di questo periodo nel Veneto passa quindi, principalmente, per la rilettura dei siti già noti. Monte Luppia e Val d'Assa sono i candidati ovvi per trovare aggiunte di età del ferro non ancora distinte stilisticamente dal bronzino circostante: una campagna RTI sistematica su pannelli già catalogati potrebbe rivelare stratificazioni non visibili a occhio nudo. Il periodo venetico, specifico del territorio veneto, ha una produzione epigrafica nota (stele funerarie, iscrizioni votive su bronzi), ma non risultano documentate iscrizioni venetiche su supporto rupestre: la loro esistenza è possibile ma non dimostrata.

% ============================================================
\section{I luoghi persistenti: quattro tipi di attrattore}
% ============================================================

L'analisi periodo per periodo rivela un dato trasversale: alcuni tipi di luogo ricorrono come attrattori di attività umana indipendentemente dall'epoca. Questi "luoghi persistenti" non sono una teoria astratta, ma un pattern documentato empiricamente in Valle Camonica, in Trentino-Alto Adige, in Scandinavia, e nelle Alpi svizzere e tirolesi. I siti rupestri tendono a concentrarsi in luoghi che soddisfano contemporaneamente criteri di più periodi, e questa sovrapposizione è di per sé un indizio di ricerca.

Il primo tipo di attrattore persistente è la sella con acqua e direttrici multiple: un passo o una forcella che mette in comunicazione due o più bacini, con sorgente a meno di ottanta metri, e da cui partono almeno due percorsi in direzioni diverse. Questo tipo di luogo è frequentato dal Mesolitico all'età del ferro senza discontinuità apparente. In Veneto, i passi delle Piccole Dolomiti (Borcola, Campogrosso) e alcuni valichi cadorini hanno questa morfologia; vanno esaminati prioritariamente.

Il secondo tipo è la parete calcarea verticale con overhang a quota intermedia: una superficie rocciosa protetta, tra quattrocento e milleduecento metri, con visibilità su una vallata. Questo tipo di luogo è prevalentemente bronzino, ma può avere aggiunte eneolitiche e dell'età del ferro. La Lessinia settentrionale ha decine di pareti di questo tipo; nessuna è stata sistematicamente cercata.

Il terzo tipo è la superficie calcarea orizzontale o subpianeggiante vicina a sentieri di quota: un pianoro o una cengia con roccia liscia, percorsa da un sentiero di mezza montagna, a quote tra ottocento e millecinquecento metri. Questo tipo di luogo ospita prevalentemente coppelle (Neolitico e Bronzo antico), ma può avere segni di tutti i periodi. I pianori carsici della Lessinia meridionale (Malga San Giorgio, altopiano di Campofontana) hanno morfologie di questo tipo.

Il quarto tipo è il crinale panoramico su due o più versanti, con morfologia a lama o a cresta arrotondata, a quote tra milleduecento e duemila metri. Questo tipo di luogo ha funzione prevalente di avvistamento e probabilmente non ospita arte rupestre in senso stretto (la roccia esposta ai venti di cresta si erode rapidamente), ma è un nodo logistico che connette altri siti. La sua importanza è metodologica: permette di identificare i percorsi tra siti, non i siti stessi.

